\section{Virasoro Algebra}

\subsection{Mode Expansion of the Stress Energy Tensor}

Starting from an infinitesimal conformal transformation in the complex coordinate parameterized by a holomorphic function
\[
    w = z + \epsilon(z),
\]
we recall that the conserved Noether charge associated with this local symmetry can be expressed as a contour integral of the stress-energy tensor \(T(z)\) around a closed path \(C\):
\[
    Q_{\epsilon} = \frac{1}{2\pi i} \oint_C \mathrm{d}z \, \epsilon(z) T(z).
\]

To systematically study these charges and the algebraic structure they form, it is extremely convenient to expand our local fields into discrete modes. On the cylinder, where the spatial coordinate is naturally compact and periodic, any generic local field \(\phi(\sigma, \tau) = \phi(w, \bar{w})\) admits a standard Fourier expansion:
\[
    \phi(w, \bar{w}) = \sum_{n, \bar{n} = -\infty}^{+\infty} \phi_{n, \bar{n}} e^{-\frac{2\pi i}{L}(n w + \bar{n} \bar{w})}.
\]
When we map this theory back to the complex plane, we must take into account the transformation law for a primary field of conformal weights \((h, \bar{h})\). Since the scaling dimensions dictate how the field rescales, this Fourier series translates directly into a Laurent expansion around the origin:
\[
    \phi(z, \bar{z}) = \sum_{n, \bar{n} = -\infty}^{+\infty} \phi_{n, \bar{n}} z^{-n-\Delta} \bar{z}^{-\bar{n}-\bar{\Delta}}.
\]
By using the Cauchy integral theorem, we can isolate and extract the individual mode operators \(\phi_{n, \bar{n}}\) via the inversion formula:
\[
    \phi_{n, \bar{n}} = \oint_0 \frac{\mathrm{d}z}{2\pi i} \oint_0 \frac{\mathrm{d}\bar{z}}{2\pi i} z^{n+\Delta-1} \bar{z}^{\bar{n}+\bar{\Delta}-1} \phi(z, \bar{z}).
\]

We now apply this general formalism specifically to the stress-energy tensor. As we established earlier, the tensor splits into independent chiral (\(T(z)\)) and anti-chiral (\(\bar{T}(\bar{z})\)) components, which possess conformal weights \((2,0)\) and \((0,2)\) respectively. Therefore, setting \(\Delta = 2\), their Laurent expansions take the special form:
\[
    T(z) = \sum_{n=-\infty}^{+\infty} L_n z^{-n-2} \;, \quad \bar{T}(\bar{z}) = \sum_{n=-\infty}^{+\infty} \bar{L}_n \bar{z}^{-n-2}.
\]
The expansion coefficients \(L_n\) and \(\bar{L}_n\) are the celebrated Virasoro modes. Inverting the series through contour integration, these modes are given explicitly by:
\[
    L_n = \oint_0 \frac{\mathrm{d}z}{2\pi i} z^{n+1} T(z) \;, \quad \bar{L}_n = \oint_0 \frac{\mathrm{d}\bar{z}}{2\pi i} \bar{z}^{n+1} \bar{T}(\bar{z}).
\]

Notice the profound connection between these modes and the conserved charges we introduced at the beginning. Because the transformations \(\epsilon(z)\) are holomorphic, they can be expanded in a basis of polynomials \(\epsilon_n(z) = -z^{n+1}\). Inserting this expansion into the integral for \(Q_{\epsilon}\), the general conserved conformal charge becomes a direct linear combination of the Virasoro modes:
\begin{equation}
    Q_{\epsilon} = \sum_{n=-\infty}^{+\infty} \epsilon_n L_n.
    \label{eq:conformal_charge_virasoro}
\end{equation}
Consequently, the holomorphic conformal generators \(\{L_n\}\) form the exact basis of a linear Lie algebra, just as the anti-holomorphic modes \(\{\bar{L}_n\}\) form an independent copy of this algebra for the anti-chiral sector.

\subsection{Virasoro Algebra from Stress Energy Tensor OPE}

We can now compute the exact commutation relations of the Virasoro modes \(L_n\). Starting from their definition via contour integrals:
\[
    L_n = \oint_0 \frac{\mathrm{d}w}{2\pi i} w^{n+1} T(w) \;, \quad L_m = \oint_0 \frac{\mathrm{d}z}{2\pi i} z^{m+1} T(z),
\]
we can express their commutator \([L_m, L_n]\) using the radial ordering formalism we developed earlier. The commutator translates into the difference of two contour integrals, which can be deformed into a single integral of the radially ordered product, where the \(z\) contour tightly encloses the point \(w\):
\[
    [L_m, L_n] = \frac{1}{(2\pi i)^2} \oint_0 \mathrm{d}w \, w^{n+1} \oint_w \mathrm{d}z \, z^{m+1} T(z)T(w).
\]

To evaluate this, we insert the singular part of the \(T(z)T(w)\) Operator Product Expansion into the integral:
\[
    [L_m, L_n] = \frac{1}{(2\pi i)^2} \oint_0 \mathrm{d}w \, w^{n+1} \oint_w \mathrm{d}z \, z^{m+1} \left[ \frac{c/2}{(z-w)^4} + \frac{2T(w)}{(z-w)^2} + \frac{\partial_w T(w)}{z-w} \right].
\]

We can evaluate the inner integral over \(z\) using Cauchy's integral formula and its derivatives:
\[
    f(w) = \oint_w \frac{\mathrm{d}z}{2\pi i} \frac{f(z)}{z-w} \;, \quad \frac{\mathrm{d}^k f(w)}{\mathrm{d}w^k} = k! \oint_w \frac{\mathrm{d}z}{2\pi i} \frac{f(z)}{(z-w)^{k+1}}.
\]
We divide the calculation into three parts, corresponding to the three poles in the OPE.

\subsubsection{Fourth-Order Pole}

The fourth-order pole gives a purely numerical contribution proportional to the central charge \(c\):
\[
    \begin{aligned}
        I_1 & = \frac{1}{2\pi i} \oint_0 \mathrm{d}w \, w^{n+1} \oint_w \frac{\mathrm{d}z}{2\pi i} z^{m+1} \frac{c/2}{(z-w)^4} \\
            & = \frac{c/2}{3!} \frac{1}{2\pi i} \oint_0 \mathrm{d}w \, w^{n+1} \partial_w^3 (w^{m+1})                          \\
            & = \frac{c}{12} (m+1)m(m-1) \frac{1}{2\pi i} \oint_0 \mathrm{d}w \, w^{m+n-1}.
    \end{aligned}
\]

Since the contour integral of \(w^k\) around the origin vanishes unless \(k = -1\) (where it equals \(2\pi i\)), this term is non-zero only when \(m+n = 0\):
\[
    I_1 = \frac{c}{12} (m^3 - m) \delta_{m+n, 0}.
\]

\subsubsection{Second-Order Pole}

The second-order pole contributes:
\[
    \begin{aligned}
        I_2 & = \frac{1}{2\pi i} \oint_0 \mathrm{d}w \, w^{n+1} \oint_w \frac{\mathrm{d}z}{2\pi i} z^{m+1} \frac{2T(w)}{(z-w)^2} \\
            & = \frac{2}{2\pi i} \oint_0 \mathrm{d}w \, w^{n+1} \partial_w (w^{m+1}) T(w)                                        \\
            & = 2(m+1) \frac{1}{2\pi i} \oint_0 \mathrm{d}w \, w^{m+n+1} T(w).
    \end{aligned}
\]

Recognizing the integral as the exact definition of the mode \(L_{m+n}\), we get:
\[
    I_2 = 2(m+1) L_{m+n}.
\]

\subsubsection{First-Order Pole}

Finally, the first-order pole yields:
\[
    \begin{aligned}
        I_3 & = \frac{1}{2\pi i} \oint_0 \mathrm{d}w \, w^{n+1} \oint_w \frac{\mathrm{d}z}{2\pi i} z^{m+1} \frac{\partial_w T(w)}{z-w} \\
            & = \frac{1}{2\pi i} \oint_0 \mathrm{d}w \, w^{m+n+2} \partial_w T(w).
    \end{aligned}
\]
Using integration by parts on the \(w\) contour, this becomes:
\[
    I_3 = - \frac{1}{2\pi i} \oint_0 \mathrm{d}w \, \partial_w (w^{m+n+2}) T(w) = - (m+n+2) L_{m+n}.
\]

\subsubsection{Virasoro Commutation Relations}

Summing the three contributions (\(I_1 + I_2 + I_3\)), we obtain the fundamental commutation relations of the modes:
\begin{equation}
    [L_m, L_n] = (m-n) L_{m+n} + \frac{c}{12} (m^3 - m) \delta_{m+n, 0}.
    \label{eq:virasoro_algebra}
\end{equation}

This equation defines the \textbf{Virasoro algebra}. It is the infinite-dimensional algebra generating the holomorphic conformal transformations in two-dimensional CFT.

By repeating the exact same procedure for the anti-holomorphic stress tensor \(\bar{T}(\bar{z})\), we find that its modes \(\bar{L}_n\) satisfy an identical algebra:
\begin{equation}
    [\bar{L}_m, \bar{L}_n] = (m-n) \bar{L}_{m+n} + \frac{c}{12} (m^3 - m) \delta_{m+n, 0}.
    \label{eq:anti-holomorphic_virasoro_algebra}
\end{equation}

Furthermore, because the holomorphic and anti-holomorphic sectors are completely independent, their respective generators commute:
\[
    [L_m, \bar{L}_n] = 0.
\]

Thus, the full local conformal symmetry in two dimensions is dynamically governed by two independent, commuting copies of the Virasoro algebra. This rich and rigidly constrained algebraic structure is the cornerstone that makes two-dimensional conformal field theories exactly solvable in many cases.

\subsection{The Full Conformal Algebra and the Classical Limit}

In other words, the complete local conformal algebra in two dimensions is given by the direct product of the holomorphic and anti-holomorphic sectors:
\[
    \Gamma = \text{Vir}_c \otimes \overline{\text{Vir}}_c
\]
where \(\text{Vir}_c\) denotes the Virasoro algebra with central charge \(c\), explicitly written as:
\begin{equation}
    \begin{aligned}
        [L_n, L_m]             & = (n-m)L_{n+m} + \frac{c}{12}n(n^2-1)\delta_{n,-m},       \\
        [L_n, \bar{L}_m]       & = 0,                                                      \\
        [\bar{L}_n, \bar{L}_m] & = (n-m)\bar{L}_{n+m} + \frac{c}{12}n(n^2-1)\delta_{n,-m}.
    \end{aligned}
    \label{eq:full_virasoro_algebra}
\end{equation}
Consequently, the physical conformal fields of the theory must lie in irreducible representations of this full algebra \(\Gamma\).

It is particularly illuminating to consider the limit where the central charge vanishes. If we set \(c = 0\), the Virasoro algebra rigorously reduces to:
\[
    [L_m, L_n] = (m-n)L_{m+n},
\]
which is precisely the \textbf{Witt algebra}, namely the classical Lie algebra of holomorphic vector fields defined on the circle.
Therefore, mathematically, the Virasoro algebra represents the unique non-trivial central extension of the Witt algebra. Physically, the central term proportional to \(c\) is a purely quantum effect, arising directly from the short-distance singularities in the operator product expansion.

This profound observation perfectly clarifies the relation between the classical discussion of conformal symmetry and its quantum counterpart:
\begin{itemize}
    \item \textbf{Classically}, the infinitesimal conformal transformations strictly form the Witt algebra without any anomalous extension;
    \item \textbf{Quantum mechanically}, the presence of the conformal anomaly forces the physical generators to satisfy the Virasoro algebra.
\end{itemize}

Historically, the profound understanding of these transformations originated in the context of string theory, where the fields defined on the two-dimensional worldsheet of the string transform precisely under the action of the Virasoro algebra. Mathematically, this constitutes a quintessential example of an infinite-dimensional Lie algebra. Its internal product is given by the mode commutator \([L_m, L_n]\), which strictly satisfies the Jacobi identity. The emergence of this algebraic structure not only revolutionized theoretical physics but also sparked an intense, dedicated mathematical effort to study and classify infinite-dimensional Lie algebras.

The overarching strategy to completely solve a two-dimensional Conformal Field Theory relies entirely on this algebraic framework. Despite the immense size of the Hilbert space and the local operator algebra, every physical state and field must fall into specific irreducible representations of the Virasoro algebra. Consequently, the full Hilbert space rigidly decomposes into a direct sum of these infinite-dimensional representations. By exploiting this representation theory, we can systematically classify the operator content of the theory and analytically compute correlation functions, which is the fundamental mechanism making two-dimensional conformal models exactly solvable.

\subsubsection{The Global Conformal Subalgebra}

An important finite-dimensional subalgebra of the full Virasoro algebra is generated by the following three modes:
\[
    L_{-1}, \quad L_0, \quad L_1.
\]
It is crucial to notice that the central extension term, \(\frac{c}{12}m(m^2-1)\delta_{m+n,0}\), strictly vanishes for \(m, n \in \{-1, 0, 1\}\). Because \(m^3 - m = 0\) on this specific subspace, the central anomaly disappears entirely. Consequently, if we restrict our algebra to these generators, the commutation relations close without any central charge contribution:
\begin{equation}
    [L_0, L_{-1}] = L_{-1} \;, \quad [L_0, L_1] = -L_1 \;, \quad [L_1, L_{-1}] = 2L_0.
    \label{eq:global_conformal_algebra}
\end{equation}

These three specific modes are precisely the generators of the \textbf{global conformal group} (the Möbius transformations). Algebraically, they form the Lie algebra \(\mathfrak{sl}(2, \mathbb{C})\) in the holomorphic sector, which is exactly what we expect from the global conformal transformations defined on the Riemann sphere.

This elegantly confirms a fundamental physical property: \uline{the central extension affects only the genuinely local, infinite-dimensional part of the conformal symmetry}. While the local conformal symmetry inevitably acquires a quantum anomaly, \textit{the global conformal symmetry remains perfectly unbroken in the quantum theory}.

\subsection{Interpretation of the Central Charge}

The central charge \(c\) is arguably the most important macroscopic parameter of a two-dimensional conformal field theory. As we established, while the classical algebra of infinitesimal conformal transformations is the Witt algebra, the quantum mechanical framework modifies this algebra by introducing a central extension. This quantum modification mathematically originates from the singular fourth-order pole in the \(TT\) Operator Product Expansion. The coefficient \(c\) precisely measures the size of this quantum anomaly.

It is extremely useful to emphasize the following fundamental points regarding \(c\):
\begin{itemize}
    \item \(c\) is a dynamical parameter and is not uniquely fixed by conformal symmetry alone.
    \item Different physical theories generally correspond to distinct values of \(c\).
    \item \(c\) actively counts, in a broad sense, the \textit{number of effective degrees of freedom of the theory}. We can plainly see this additive property from our earlier concrete examples:
          \begin{itemize}
              \item a single free real scalar boson has \(c = 1\),
              \item a single free Majorana fermion has \(c = 1/2\).
          \end{itemize}
\end{itemize}

This intuitive statement (that the central charge counts the effective degrees of freedom) can be made mathematically rigorous with advanced methods that go beyond the immediate scope of these lectures (most notably through Zamolodchikov's famous \(c\)-theorem, which proves that a generalized function \(c\) strictly decreases along Renormalization Group flows, formalizing the physical loss of high-energy degrees of freedom as a system cools down).

Ultimately, the physical significance of the central charge is multifold and pervasive across the whole theory:
\begin{itemize}
    \item It formally measures the \textbf{conformal anomaly} (the quantum breaking of classical scale invariance when placing the theory on a curved background).
    \item It actively appears in the finite-size geometric dependence of the vacuum energy on the cylinder (generating the \textbf{Casimir effect}).
    \item It universally controls the asymptotic growth of the density of states at extremely high energies through \textbf{Cardy's formula}, effectively determining the high-energy thermodynamic entropy of the system.
    \item It plays a central, rigid role in the exact classification of the so-called \textbf{minimal models} (rational CFTs with a finite number of primary fields).
\end{itemize}

Thus, the precise numerical value of \(c\) is not merely an algebraic curiosity, but a quantity that encapsulates the deepest dynamical and thermodynamic information about the physical theory.

\subsection{Action of Virasoro Generators on Primary Fields}

We now want to determine the explicit action of the Virasoro generators on a local primary field. To do this, we must compute the commutator \([L_n, \Phi(w, \bar{w})]\). In the framework of radial quantization, we must be incredibly careful with the ordering of operators: the commutator translates geometrically into the difference of two contour integrals over \(z\), one where \(|z| > |w|\) and one where \(|z| < |w|\). Applying the exact same contour deformation procedure we used for the \(TT\) commutator, this difference collapses into a single contour integral over \(z\) tightly wrapping the insertion point \(w\):
\[
    [L_n, \Phi(w, \bar{w})] = \oint_w \frac{\mathrm{d}z}{2\pi i} z^{n+1} T(z)\Phi(w, \bar{w}).
\]
We can now evaluate this integral by starting from the defining Operator Product Expansion of the stress-energy tensor with a primary field:
\[
    T(z)\Phi(w, \bar{w}) \sim \frac{h\Phi(w, \bar{w})}{(z-w)^2} + \frac{\partial_w \Phi(w, \bar{w})}{z-w}.
\]
Inserting the singular part of this OPE directly into our contour integral, we obtain:
\[
    [L_n, \Phi(w, \bar{w})] = \oint_w \frac{\mathrm{d}z}{2\pi i} z^{n+1} \left[ \frac{h\Phi(w, \bar{w})}{(z-w)^2} + \frac{\partial_w \Phi(w, \bar{w})}{z-w} \right].
\]
By applying Cauchy's integral formula and its first derivative formula to the analytic function \(f(z) = z^{n+1}\) evaluated at the pole \(z = w\), we can easily extract the residues for both terms:
\begin{equation}
    [L_n, \Phi(w, \bar{w})] = (w^{n+1}\partial_w + h(n+1)w^n)\Phi(w, \bar{w}).
    \label{eq:virasoro_primary_action}
\end{equation}
This formula is extremely important and represents a key result in two-dimensional conformal field theory. It shows explicitly how the infinite set of local Virasoro generators act dynamically on primary fields as differential operators.

Furthermore, we can observe the consistency of this result with the global subgroup. If we restrict this action to the modes \(n = -1, 0, 1\), the differential operators accurately reproduce the classical action of the globally defined transformations. Specifically, \(L_{-1}\) generates spatial translations (\(\partial_w\)), \(L_0\) generates dilatations and rotations (\(w\partial_w + h\)), and \(L_1\) generates special conformal transformations (\(w^2\partial_w + 2hw\)).